\setcounter{chapter}{0}
\renewcommand{\thechapter}{1}
\chapter{The Integers}\label{ch:1}
\setcounter{equation}{0}	        % To start with Equation 1
\counterwithin{equation}{section}	% Equation will be will be 1.1 1.2 1.3   ... 

\section{Terminology of Sets}\label{sec:1.1}	% NO  exercise
	NO Exercise

\newpage		% Forcing a page break

\section{Basic Properties}\label{sec:1.2}	
\subsection*{Exercises}\label{sec:1.2_Exercises}

\subsubsection{Ex 1.2 (1)}\label{sec:Ex_1.2.1}
% Ex 1 statement

If $n, m$ are integers $\geq 1$ and $n \geq m$, define the \tbf{binomial coefficients}
$$ \binom{n}{k} = \frac{n!}{m! (n-m)!} $$
As usual, $n! = n \cdot (n-1) \cdots 1$ is the product of the first $n$ integers. We define $0! = 1$ and 
$\binom{n}{0} = 1$. Prove that 
$$
\binom{n}{m-1} + \binom{n}{m} = \binom{n+1}{m}\,.
$$

% Ex 1 solution
%
\paragraph{Solution:}{

\begin{align*}
    \binom{n}{m-1} + \binom{n}{m} &= \frac{n!}{(m-1)!(n-m+1)!} +  \frac{n!}{m!(n-m)!}\\
	&= \frac{n!}{ (m-1)! (n-m+1)(n-m)!} + \frac{n!}{m (m-1)!(n-m)!} \\  
	&= \frac{m \cdot n!  +   (n - m + 1) \cdot n!}{ m \cdot (m-1)! (n-m+1)(n-m)!} \\  
	&= \frac{(n+1) n!}{ m \cdot (m-1)! (n-m+1)(n-m)!} \\  	        
    &= \frac{(n+1)!}{ m! (n+1-m)!} = \binom{n+1}{m}\,.
\end{align*}

}

\subsubsection{Ex 1.2 (2)}\label{sec:Ex_1.2.2}
% Ex 2 statement
Prove by induction that for any integers $x, y$ we have
$$ 
(x + y)^n = \sum_{i=0}^{n} \binom{n}{i} x^i y^{n-i} = y^n + \binom{n}{1} x y^{n-1} + \binom{n}{2} x^2 y^{n-2}
+ \cdots + x^n\,.
$$

%
% Ex 2 solution
%
\paragraph{Solution:}{
According to the first form of induction, we must take the above statement as a proposition $P(n)$ with variable $n \in \N$, and we must
\begin{enumerate}[a)]
\item Prove $P(1)$;
\item Prove $P(n) \implies P(n+1)$. 
\end{enumerate}

Let now proceed:

\begin{enumerate}[A)]
\item Proof of $P(1)\;$:
\begin{align*}
	\sum_{i=0}^{1} \binom{1}{i} x^i y^{n-i} &= \binom{1}{0} x^0 y^1 + \binom{1}{1} x^1 y^{1-1}\\
	                                        &= 1 \cdot 1 \cdot y + 1 \cdot x \cdot y^0\\
	                                        &= y + x \qed
\end{align*}
\item Proof that $P(n) \implies P(n+1)\;$:
This amounts to prove the identity
$$ (x + y)^{n+1} = (x + y) \sum_{i=0}^{n} \binom{n}{i} x^i y^{n-i} = \sum_{i=0}^{(n+1)} \binom{n+1}{i} x^i y^{n-i+1}\,,$$
likely exploiting the identity proved in exercise 1 above, namely
$$\binom{n}{i-1} + \binom{n}{i} = \binom{n+1}{i}\,.$$
 
The proof starts here :
\begin{align*}
P &=  (x + y)^{n+1} = (x + y) \cdot (x + y)^n\\
  &=  x \cdot (x + y)^n + y \cdot (y + x)^n \\
  &=  \binom{n}{0} x^1 y^{n} + \binom{n}{1} x^2 y^{n-1} + \cdots + \binom{n}{n-1} x^{n} y^{1}  + \binom{n}{n} x^{n+1} y^{0}\\
  &+  \binom{n}{0} y^1 x^{n} + \binom{n}{1} y^2 x^{n-1} + \cdots + \binom{n}{n-1} y^{n} x^{1}  + \binom{n}{n} y^{n+1} x^{0}
\end{align*}
 
We now make use of the trivial identity 
$$\binom{n}{n-k} = \binom{n}{k}\,,$$ 
and rewrite the above expression for $(x + y)^{n+1}$ in a form that helps identifying the \tbi{two} binomial coefficients which multiply the \tbi{same term} $x^j y^k$ in the above expansion:
\begin{align*}
P &=  (x + y)^{n+1} = (x + y) \cdot (x + y)^n  = x (x + y)^n + y (x + y)^n\\
  &=  \binom{n}{0} x^1 y^{n} + \binom{n}{1} x^2 y^{n-1} + \cdots + \binom{n}{1} x^{n} y^{1}  + \binom{n}{0} x^{n+1} y^{0}\\
  &+  \binom{n}{0} y^1 x^{n} + \binom{n}{1} y^2 x^{n-1} + \cdots + \binom{n}{1} y^{n} x^{1}  + \binom{n}{0} y^{n+1} x^{0}\,.\\
  &=  x^{n+1} + y^{n+1}\\
  &+  \binom{n}{0} x^1 y^{n} + \binom{n}{1} x^2 y^{n-1} + \cdots + \binom{n}{2} x^{n-1} y^{2} + \binom{n}{1} x^{n} y^{1}\\
  &+  \binom{n}{0} y^1 x^{n} + \binom{n}{1} y^2 x^{n-1} + \cdots + \binom{n}{2} y^{n-1} x^{2} + \binom{n}{1} y^{n} x^{1}
\end{align*}

After a little re-arrangement:  
\begin{align*}
P &=  x^{n+1} + y^{n+1}\\
  &+  \binom{n}{0} x^1 y^{n} + \binom{n}{1} x^2 y^{n-1} + \cdots + \binom{n}{2} x^{n-1} y^{2} + \binom{n}{1} x^{n} y^{1} \\
  &+  \binom{n}{0} y^1 x^{n} + \binom{n}{1} y^2 x^{n-1} + \cdots + \binom{n}{2} y^{n-1} x^{2} + \binom{n}{1} y^{n} x^{1}\\
  &=  x^{n+1} + y^{n+1} \\
  &+  \left[ \binom{n}{0} +  \binom{n}{1} \right] \cdot (x^{1} y^{n} + y^1 x^n)\\
  &+  \left[ \binom{n}{1} + \binom{n}{2} \right] \cdot (x^{2} y^{n-1} + y^2 x^{n-1})\\
  &+  \cdots \text{hence, using the result of Exercise 1:}\\
P &=  x^{n+1} + y^{n+1} + \binom{n+1}{1} (x^{1} y^{n} + y^1 x^n) + \binom{n+1}{2} (x^{2} y^{n-1} + y^2 x^{n-1}) + \cdots \\
  &=  \sum_{i=0}^{n+1} \binom{n+1}{i} x^i y^{n+1-i} \qed
\end{align*}

\end{enumerate}

\subsubsection{Ex 1.2 (3)}\label{sec:Ex_1.2.3}
% Ex 3 statement
Prove the following statements for all positive integers:
\begin{enumerate}[(a)]
\item $1 + 3 + 5 + \cdots + (2n - 1) = n^2$
\item $1^2 + 2^2 + \cdots + n^2 = n (n+1)(2n + 1)/6$
\item $1^3 + 2^3 + \cdots + n^3 = [n (n+1)/2]^2$
\end{enumerate}

%
% Ex 3 solution
%
\paragraph{Solution:}{
\begin{enumerate}[(a)]
\item $\sum_{i=1}^{n} (2n-1) = 2 \sum_{i=1}^{n} i - \sum_{i=1}^{n} 1 = 2 \frac{n(n+1)}{2} - n = n(n+1) - n = n^2\,.$ 
\item The statement is true for $n = 1\,: 1^2 = 1 \cdot 2 \cdot 3/6 = 6/6$. By induction, it is enough to prove that \tbi{if} the sum of the squares $i^2$ of the first $n$ integers equals $n (n+1)(2n + 1)/6$ \tbi{then} by adding $(n+1)^2$ to this sum one obtains $m (m+1)(2m + 1)/6$, with $m=n+1$. Indeed both expressions reduce to $2 n^3 + 9 n^2 + 13 n + 6$, hence the identity is proved.
\item The statement is true for $n = 1\,: 1^3 = [1\cdot(1+1)/2)]^2 = 1$. By induction, it is enough to prove that \tbi{if} the sum of the cubes $i^3$ of the first $n$ integers equals $[n (n+1)/2]^2$ \tbi{then} by adding $(n+1)^3$ to this sum one obtains $[m (m+1)/2]^2$, with $m=n+1$. Indeed both expressions reduce to $n^4+ 6 n^3 + 13 n^2 + 12 n + 4$, hence the identity is proved.
\end{enumerate}
}

\subsubsection{Ex 1.2 (4)}\label{sec:Ex_1.2.4}
% Ex 4 statement
Prove that 
$$ \left( 1 + \frac{1}{1} \right)^1 \left( 1 + \frac{1}{2} \right)^2 \cdots {\left( 1 + \frac{1}{n-1} \right)}^{(n-1)} = \frac{n^{(n-1)}}{(n-1)!} $$
%
% Ex 4 solution
%
\paragraph{Solution:}{
Let $A(n)$ and $B(n)$ be the left and right sides of the above equation, 
whose identity\footnote{Note $n^{n-1}/(n-1)!$ \tit{numerically equals} $n^n / n!$.} 
we want to prove:
\begin{align*}
A(n) &= \prod_{i=1}^{n-1} \left( 1 + \frac{1}{i} \right)^i\,, &B(n) &= \frac{n^{n-1}}{(n-1)!} = \frac{n^n}{n!}\,.
\end{align*}

The \tbi{identity holds\footnote{$2$ is the lowest value of $n$ for which the equation $A(n) = B(n)$ was declared to hold.} for}  $n=2\;$: 
\begin{align*}
A(2) &= \left( 1 + \frac{1}{1} \right)^1\,= 2 &B(2) &= \frac{2^{2-1}}{(2-1)!} \;= \frac{2^2}{2!} \;= 2 \,.
\end{align*}

If true, the identity $A(n) = B(n)$ for all $n \geq 2$ should follow by\\ \tbi{induction}, i.e. by proving that 
$$ A(n) = B(n) \text{\:}\implies \text{\:}  A(n+1) = B(n+1) \:\: \forall \:\: n \geq 2$$

Indeed:

\begin{align*}
A(n+1) &= \prod_{i=1}^n \left( 1 + \frac{1}{i} \right)^i\\
       &= A(n) \cdot  \left( 1 + \frac{1}{n} \right)^n\\
       &= B(n) \cdot  \left( 1 + \frac{1}{n} \right)^n\\
       &= \frac{n^{n-1}}{(n-1)!}  \cdot  \left(\frac{n+1}{n}\right)^n = (n+1)^n \cdot \left(\frac{n^{n-1}}{n^n}\right) \cdot \frac{1}{(n-1!)}\\
       &= (n+1)^n \cdot \frac{1}{n} \cdot \frac{1}{(n-1!)}\\
       &= \frac{(n+1)^n}{n!} = B(n+1)\,. \qed   
\end{align*}

}


\subsubsection{Ex 1.2 (5)}\label{sec:Ex_1.2.5}
% Ex 5 statement

Let $x$ be a real number. Prove that there exists an integer $q$ and a real number $s$
with $0 \leq s < 1$ such that $x = q + s$, and that $q, s$ are uniquely determined. 
Can you deduce the Euclidean algorithm from this result without using induction?\footnote{Search \tit{Euclidean algorithm}  on Wikipedia.}

\paragraph{Solution:}{
For any real number $x$, there exist a unique integer $q$ and a unique real number $s \geq 0$
such that $x = q + s$ \tbi{if and only if} we additionally enforce the constraint $s < 1$
\footnote{Without the bound $s < 1$, uniqueness is lost: $5.5 = 5 + 0.5$ and $5.5 = 4 + 1.5$.}.
Assuming this standard definition of the integer and fractional part decomposition, here is the proof divided into
existence and uniqueness, where the \tit{floor} function $\lfloor x \rfloor$ denotes the \tit{integer part} of a real number $x$. 
\begin{enumerate}
\item \tbf{Existence of $q$ and $s$}\\
$q = \lfloor x \rfloor \leq x$, and $ s = x - \lfloor x \rfloor < 1$

\item \tbf{Uniqueness of $q$ and $s$}\\
Assume there are two different pairs of representations for the same real number $x$:\\
$x = q_1 + s_1$ and $x = q_2 + s_2$, where $$q_1, q_2 \in \Z \text{ and } 0 \leq s_1, s_2 < 1\,.$$
Equating both expressions for $x$ gives: 
$$q_1 - q_2 =  s_2 - s_1$$
Since the difference between two integers $(q_1 - q_2)$ must be an integer, the term $(s_2 - s_1)$ must also evaluate to an integer.
Given $0 \leq s_2 < 1$ and $-1 < -s_1 \leq 0$, adding these inequalities yields: $-1 < s_2 - s_1 < 1$. The only integer that strictly
lies between $-1$ and $1$ is $0$. Therefore: $s_2 - s_1 = 0 \rightarrow s_1 = s_2$\\
Substituting $s_2 - s_1 = 0$ back into our rearranged equation:
$$q_1 - q_2 = 0 \rightarrow q_1 = q_2$$
Because both $q_1 = q_2$ and $s_1 = s_2$, the representation is uniquely determined $\qed$.
\end{enumerate}

\subsubsection{Ex 1.2 (6) -- Euclidean algorithm}\label{sec:Ex_1.2.6}

The existence and uniqueness of the decomposition of any real number into an integer (quotient) and a fractional (remainder) part—is formally known as the Quotient-Remainder Theorem or Euclidean Division. 
For any two integers $a$ and $b$ (where $b \neq 0$), there exists a unique quotient $q$ 
and remainder $r$ such that: $a = b q + r$ where $0 \leq r < \abs{b}$.\\

From this decomposition, a crucial mathematical property arises:
\begin{itemize}
\item Any common divisor $cd$ that perfectly divides both $a$ and $b$ will \tbf{also perfectly divide their difference $|b - a|$}, because
the integers $cd, x, y$ are such that $a = cd*x$; $b = cd*y$, hence $|b - a| = cd * |x - y|$. 
\item Finding the GCD of the two numbers can thus be done by iteratively replacing the largest one, at each step, with their difference, until they become equal to the GCD of $a$ and $b$.
\end{itemize}

\newpage		% Forcing a page break

\section{Greatest Common Divisor}\label{sec:1.3}
% NO EXERCISE


\newpage		% Forcing a page break

\section{Unique Factorization}\label{sec:1.4}
\subsection*{Exercises}\label{sec:1.4_Exercises}

\subsubsection{Ex 1.4 (1)}\label{sec:Ex_1.4.1}
% Ex 1 statement

\subsubsection{Ex 1.4 (2)}\label{sec:Ex_1.4.2}
% Ex 2 statement

\subsubsection{Ex 1.4 (3)}\label{sec:Ex_1.4.3}
% Ex 3 statement

\subsubsection{Ex 1.4 (4)}\label{sec:Ex_1.4.4}
% Ex 4 statement

\newpage		% Forcing a page break

\section{Equivalence Relations and Congruences}\label{sec:1.5}
\subsection{Exercises}\label{sec:1.5_Exercises}

\subsubsection{Ex 1.5 (1)}\label{sec:Ex_1.5.1}
% Ex 1 statement

\subsubsection{Ex 1.5 (2)}\label{sec:Ex_1.5.2}
% Ex 2 statement

\subsubsection{Ex 1.5 (3)}\label{sec:Ex_1.5.3}
% Ex 3 statement

\subsubsection{Ex 1.5 (4)}\label{sec:Ex_1.5.4}
% Ex 4 statement

\subsubsection{Ex 1.5 (5)}\label{sec:Ex_1.5.5}
% Ex 5 statement

\subsubsection{Ex 1.5 (6)}\label{sec:Ex_1.5.6}
% Ex 6 statement

\subsubsection{Ex 1.5 (7)}\label{sec:Ex_1.5.7}
% Ex 7 statement

\subsubsection{Ex 1.5 (8)}\label{sec:Ex_1.5.8}
% Ex 8 statement

\subsubsection{Ex 1.5 (9)}\label{sec:Ex_1.5.9}
% Ex 9 statement

\subsubsection{Ex 1.5 (10)}\label{sec:Ex_1.5.10}
% Ex 10 statement

\newpage		% Forcing a page break

\subsubsection{Ex 1.5 (11)}\label{sec:Ex_1.5.11}
% Ex 11 statement

\subsubsection{Ex 1.5 (12)}\label{sec:Ex_1.5.12}
% Ex 12 statement




